\documentclass[11pt]{amsart}
\usepackage{amsmath,amssymb,amsthm}
\usepackage[margin=1.15in]{geometry}
\usepackage{booktabs}
\usepackage[hidelinks]{hyperref}

\newtheorem{theorem}{Theorem}
\newtheorem{corollary}[theorem]{Corollary}
\newtheorem{lemma}[theorem]{Lemma}
\newtheorem{proposition}[theorem]{Proposition}
\theoremstyle{remark}
\newtheorem{remark}[theorem]{Remark}

\newcommand{\msum}{\mathsf{msum}}
\newcommand{\disc}{\mathsf{disc}}
\newcommand{\Z}{\mathbb{Z}}

\title{Small maximal $k$-consecutive sums on a circle\\ and the clock-face problem}
\author{Victor Nidal Bassem Nagib}
\date{October 30, 2026; revised November 10, 2026}

\begin{document}

\begin{abstract}
Arrange $1,2,\dots,n$ on a circle and let $\msum(n,k)$ be the smallest possible value of the largest sum of $k$
consecutive entries, minus the average $k(n+1)/2$ of these sums. For odd $k\ge 3$ and odd $m$, Higashitani and
Kurimoto proved $\msum(mk,k)=2$ for $m\ge 2k+3$, while for $5\le m\le 2k+1$ it was only known that
$\msum(mk,k)\in\{1,2\}$, apart from two cases with $k=3$ and one case with $k=7$ settled by computer. We show that $\msum(mk,k)=1$ for $3\le m\le k+2$ and $\msum(mk,k)=2$ for
$m\ge k+4$, which settles this case. The lower bound comes from a short total-variation argument giving, for
odd $k$, $\msum(n,k)\ge 2$ for every odd $n>(k+3)\gcd(n,k)$; the upper bound is an explicit construction. For
$k=3$ the argument gives a short proof, with no computer verification, that the clock-face numbers (OEIS A066385)
are at least $\lceil 3n/2\rceil+3$ for all $n\ge 4$ except $n=5,6,9,15$. With the known matching upper bound this
gives $a(n)=\lceil 3n/2\rceil+3$ apart from these four exceptions, as conjectured by J.~E.~Schoenfield. This also
follows from the work of Higashitani and Kurimoto, except for the case $n=21$, which they credit to a personal
communication and which earlier computer searches had also settled.
\end{abstract}

\maketitle

\section{Introduction}

Let $n>k\ge 1$. For a permutation $\pi=(\pi_0,\dots,\pi_{n-1})$ of $[n]=\{1,\dots,n\}$, read cyclically
($\pi_{i+n}=\pi_i$), let
\[
  s_i=s_i(\pi)=\pi_i+\pi_{i+1}+\dots+\pi_{i+k-1}\qquad(i\in\Z_n)
\]
be its $k$-consecutive sums. They have average $k(n+1)/2$. Following Higashitani and Kurimoto~\cite{HK} put
\[
  \msum(\pi,k)=\max_{i}s_i(\pi)-\frac{k(n+1)}{2},\qquad \msum(n,k)=\min_{\pi}\msum(\pi,k).
\]
Since $\max_i s_i$ is an integer at least the average, $\msum(n,k)$ is a non-negative integer when $k(n+1)$ is even
and lies in $\frac12+\Z$ when $k(n+1)$ is odd. Anstee, Ferguson and Griggs~\cite{AFG} studied the two-sided
version $\disc(n,k)=\min_\pi\max_i|s_i(\pi)-k(n+1)/2|\ge\msum(n,k)$.

The case $k=3$ is the \emph{clock-face problem}: what is the least possible value $a(n)$ of the largest sum of
three neighbours when $1,\dots,n$ are placed around a circle? It goes back to Clark~\cite{Clark} and to a 1986
column of Martin Gardner~\cite{Gardner}, and was studied as a backtracking problem by Rolfe~\cite{Rolfe,RP}. A
version with $n=12$ was a problem of the Bundeswettbewerb Mathematik 2002~\cite{BWM}: for twelve sticks of lengths
$1,\dots,12$ in a circle, show that some three neighbours have total length at least $20$. The values are OEIS sequence A066385~\cite{OEIS}, and
$a(n)=\frac{3(n+1)}{2}+\msum(n,3)$. In 2019 J.~E.~Schoenfield observed there that $a(n)=f(n):=\lceil 3n/2\rceil+3$
for all $4\le n\le 38$ except $n=5,6,9,15$, gave an explicit arrangement showing $a(n)\le f(n)$ for all $n$,
noted that $a(n)=f(n)$ is not hard to prove for even $n>6$, and conjectured $a(n)=f(n)$ for all odd $n>15$.

Higashitani and Kurimoto~\cite{HK} determined $\msum(n,k)$ in many cases. Their Corollary~1.6 gives
$\msum(n,3)=2$ for all odd $n\ge7$ other than $9$ and $15$, which contains Schoenfield's conjecture, except that the value
$\msum(21,3)=2$ is credited there to a personal communication of T.~Yamashita~\cite[Remark~1.7]{HK}. For odd $k$
and $n=mk$ with $m$ odd they prove $\msum(mk,k)=2$ for $m\ge 2k+3$~\cite[Theorem~1.1]{HK}, and $\msum(3k,k)=1$
follows from $\disc(3k,k)=1$~\cite{AFG}. Since $1\le\msum(mk,k)\le\disc(mk,k)\le 2$~\cite{AFG}, the open cases were
odd $m$ with $5\le m\le 2k+1$, except that for $k=3$ the values $\msum(15,3)=1$ and $\msum(21,3)=2$ were
known~\cite[Remark~1.7]{HK}, and $\msum(35,7)=1$ follows from the value $\disc(35,7)=1$ computed by integer
programming by Bogdanovi\'c, Maksimovi\'c, Simi\'c and Milo\v{s}evi\'c~\cite{BMSM}.

Our first result is a lower bound for odd $k$ and all $n$.

\begin{theorem}\label{thm:lower}
Let $k\ge3$ be odd, $n>k$ and $g=\gcd(n,k)$.
\begin{enumerate}
\item[(a)] If $n$ is odd and $n>(k+3)g$, then $\msum(n,k)\ge 2$.
\item[(b)] If $n$ is even and $n\ne 2k$, then $\msum(n,k)\ge 3/2$.
\end{enumerate}
\end{theorem}

Part (a) for $n=mk$ says $\msum(mk,k)\ge2$ for odd $m\ge k+4$, improving the range $m\ge 2k+3$
of~\cite[Lemma~5.3]{HK}. Part (b) is included for completeness; much of it is contained in~\cite{HK} (for $k=3$
it is part of~\cite[Corollary~1.6(a)]{HK}), and $n=2k$ is a genuine exception, since
$\msum(2k,k)=\disc(2k,k)=\frac12$~\cite{AFG}. Our second result shows that the bound in (a) is sharp when
$k\mid n$.

\begin{theorem}\label{thm:construction}
Let $k$ and $m$ be odd integers with $k\ge1$ and $3\le m\le k+2$. Then there is a permutation of $[mk]$ all of whose
$k$-consecutive sums are at most $\frac{k(mk+1)}{2}+1$. In other words $\msum(mk,k)\le1$.
\end{theorem}

\begin{corollary}\label{cor:mk}
Let $k$ and $m$ be odd integers with $k\ge3$ and $m\ge3$. Then
\[
  \msum(mk,k)=\begin{cases}1,& m\le k+2,\\ 2,& m\ge k+4.\end{cases}
\]
\end{corollary}

For instance $\msum(n,5)=1$ for $n=15,25,35$ and $\msum(n,5)=2$ for $n=45,55,65,\dots$; previously the values for
$n=25,35,45,55$ were open. For $k=3$, together with the known upper bound $a(n)\le\lceil 3n/2\rceil+3$
\cite{AFG,OEIS}, we obtain the complete answer to the clock-face problem; our proof of the lower bound does not
depend on any computation.

\begin{corollary}\label{cor:clock}
For all $n\ge4$,
\[
  a(n)=\Bigl\lceil\frac{3n}{2}\Bigr\rceil+3,
\]
except $a(5)=10$, $a(6)=11$, $a(9)=16$ and $a(15)=25$, which are one less. (Also $a(3)=6$.)
\end{corollary}

The four exceptions are exactly the values of $n$ that the inequalities of Theorem~\ref{thm:lower} leave open
for $k=3$ (namely the odd $n\le 6\gcd(n,3)$, and $n=6$).

\section{Slack and total variation}

Throughout this section $k\ge1$, $n>k$, and $\pi$ is a permutation of $[n]$ whose $k$-consecutive sums satisfy
$s_i\le L$ for some integer $L$. Put
\[
  \sigma_i=L-s_i\ge0,\qquad T=\sum_{i\in\Z_n}\sigma_i=nL-\frac{kn(n+1)}{2},
\]
using $\sum_i s_i=k\sum_i\pi_i$. Let
\begin{gather*}
  d_i=\sigma_i-\sigma_{i+1}=s_{i+1}-s_i=\pi_{i+k}-\pi_i ,\\
  M=\sum_{i}\min(\sigma_i,\sigma_{i+1}),\qquad P=\#\{i:\sigma_i>0\text{ and }\sigma_{i+1}>0\}.
\end{gather*}

\begin{lemma}\label{lem:basic}
\begin{enumerate}
\item[(i)] $d_i\ne0$ for every $i$; in particular $\sigma_i$ and $\sigma_{i+1}$ are never both $0$.
\item[(ii)] $\sum_i|d_i|=2T-2M$.
\item[(iii)] $M\ge P$.
\end{enumerate}
\end{lemma}

\begin{proof}
(i) As $0<k<n$, the positions $i$ and $i+k$ differ, so $\pi_{i+k}\ne\pi_i$. (ii) Use
$|x-y|=x+y-2\min(x,y)$ and sum over $i$. (iii) Every pair counted by $P$ contributes at least $1$ to $M$, and every
term of $M$ is non-negative.
\end{proof}

\emph{Chains.} The map $i\mapsto i+k$ splits $\Z_n$ into $g=\gcd(n,k)$ cycles, the \emph{chains}, each of length
$\ell=n/g\ge2$. Along a chain $i_0,i_0+k,i_0+2k,\dots$ the values $v_j=\pi_{i_0+jk}$ $(j\in\Z_\ell)$ are $\ell$ distinct
integers and $v_{j+1}-v_j=d_{i_0+jk}$. Call $v_j$ a \emph{peak} if $v_j>v_{j\pm1}$ and a \emph{valley} if
$v_j<v_{j\pm1}$; going around the chain, peaks and valleys alternate, so a chain with $r\ge1$ peaks has exactly $r$
valleys. Put
\[
  D=\#\{i\in\Z_n: d_id_{i+k}<0\}.
\]
Since $i$ and $i+k$ lie in the same chain, $d_id_{i+k}<0$ exactly when $\pi_{i+k}$ is a peak or a valley of its
chain. Hence $D=\sum_c 2r_c$, where $r_c\ge1$ is the number of peaks of chain $c$; in particular $D\ge 2g$.

\begin{lemma}\label{lem:tv}
$\sum_i|d_i|\ge 2n-4g+D$.
\end{lemma}

\begin{proof}
Fix a chain with values $v_j$ and $r$ peaks. Write its peaks and valleys in cyclic order as
$V_1,P_1,V_2,P_2,\dots,V_r,P_r$. Its total variation is
\[
  \sum_j|v_{j+1}-v_j|=\sum_{t}(P_t-V_t)+\sum_t(P_t-V_{t+1})=2\sum_tP_t-2\sum_tV_t .
\]
Remove from the cyclic sequence $V_1,P_1,\dots,V_r,P_r$ the largest value $P^*$ (a peak) and the smallest value
$V^*$ (a valley). What remains is at most two paths (none if $r=1$) that alternate between peaks and valleys and have an even number
of terms (each path runs from a neighbour of one removed term to a neighbour of the other, and these neighbours are
of opposite types). So the remaining $r-1$ peaks can be matched with the remaining $r-1$ valleys so that matched
terms are adjacent in the cyclic sequence; a peak exceeds its neighbouring valleys, so each matched pair has
$P-V\ge1$. Since the $\ell$ values are distinct integers, $P^*-V^*\ge \ell-1$. Therefore the total variation is at
least $2(\ell-1)+2(r-1)$. Summing over the $g$ chains, which partition the indices $i$, gives
$\sum_i|d_i|\ge 2g(\ell-1)+2\sum_c(r_c-1)=2n-4g+D$.
\end{proof}

\begin{proposition}\label{prop:key}
$2M+D\le 2(T-n)+4g$.
\end{proposition}

\begin{proof}
Combine Lemma~\ref{lem:basic}(ii) and Lemma~\ref{lem:tv}: $2T-2M\ge 2n-4g+D$.
\end{proof}

\begin{lemma}\label{lem:pairs}
If $k$ is odd, then $n-D\le(k+1)P$.
\end{lemma}

\begin{proof}
Let $i$ be such that $d_id_{i+k}>0$; by Lemma~\ref{lem:basic}(i) these are the $n-D$ indices not counted by $D$.
Suppose that none of the $k+1$ consecutive pairs $(\sigma_{i+j},\sigma_{i+j+1})$, $0\le j\le k$, consists of two
positive numbers. No such pair consists of two zeros either (Lemma~\ref{lem:basic}(i)), so the numbers
$\sigma_i,\sigma_{i+1},\dots,\sigma_{i+k+1}$ alternate between zero and positive. As $k$ is odd, $\sigma_{i+k}$ is
zero exactly when $\sigma_i$ is positive. If $\sigma_i>0$, then $\sigma_{i+1}=0$, $\sigma_{i+k}=0$ and
$\sigma_{i+k+1}>0$, so $d_i>0>d_{i+k}$; if $\sigma_i=0$, then likewise $d_i<0<d_{i+k}$. Both contradict
$d_id_{i+k}>0$. Hence for each such $i$ there is a $j\in\{i,i+1,\dots,i+k\}$ with $\sigma_j>0$ and $\sigma_{j+1}>0$.
A given $j$ arises in this way from at most $k+1$ indices $i$ (namely $i\in\{j-k,\dots,j\}$), so $n-D\le(k+1)P$.
\end{proof}

\section{Proof of Theorem~\ref{thm:lower}}

Let $k\ge3$ be odd and $n>k$.

(a) Let $n$ be odd, so $k(n+1)/2$ is an integer, and suppose some $\pi$ has $\msum(\pi,k)\le1$. Take
$L=\frac{k(n+1)}{2}+1$; then $T=n$, and Proposition~\ref{prop:key} reads $2M+D\le 4g$. As $D\ge 2g$, we get $M\le g$
and hence $P\le g$ by Lemma~\ref{lem:basic}(iii). By Lemma~\ref{lem:pairs}, $D\ge n-(k+1)P$, so
\[
  n-(k+1)P+2P\le D+2M\le 4g ,\qquad\text{i.e.}\qquad n\le 4g+(k-1)P\le (k+3)g .
\]
So $\msum(n,k)\le 1$ forces $n\le(k+3)g$. Since $\msum(n,k)$ is an integer here, $\msum(n,k)\ge2$ whenever
$n>(k+3)g$.

(b) Let $n$ be even, so $\msum(n,k)\in\frac12+\Z$, and suppose some $\pi$ has $\msum(\pi,k)\le\frac12$. Take
$L=\frac{k(n+1)+1}{2}$; then $T=n/2$ and Proposition~\ref{prop:key} gives $2M+D\le 4g-n$. With $M\ge0$ and
$D\ge2g$ this yields $n\le 2g$. Since $g$ is odd and divides the even number $n$, $2g$ divides $n$, so $n=2g$; and
$2g=n>k$ with $g\mid k$ forces $g=k$. Hence $n=2k$. \qed

\begin{remark}\label{rem:trivial}
For odd $n$ we also have $\msum(n,k)\ge1$: if $\msum(\pi,k)=0$, take $L=k(n+1)/2$; then $T=0$, all $\sigma_i=0$ and
all $d_i=0$, contradicting Lemma~\ref{lem:basic}(i).
\end{remark}

\section{The construction}

Let $k,m\ge1$. For an array $A=(A(q,c))$ with rows $q\in\Z_m$ and columns $c\in\{0,\dots,k-1\}$ define, for
$0\le r\le k$, the \emph{staircase sums}
\[
  W_A(q,r)=\sum_{c<r}A(q+1,c)+\sum_{c\ge r}A(q,c).
\]
(So $W_A(q,0)$ is the sum of row $q$ and $W_A(q,k)=W_A(q+1,0)$.) Call $A$ \emph{good} if its columns are
permutations of $[m]$ and $W_A(q,r)\le\frac{k(m+1)}{2}+1$ for all $q$ and $r$.

\begin{lemma}\label{lem:blocks}
Let $A$ be a good $m\times k$ array and $\beta$ a permutation of $\{0,\dots,k-1\}$. Then
$\pi_{qk+c}=m\beta(c)+A(q,c)$ $(0\le q<m,\ 0\le c<k)$ defines a permutation of $[mk]$ all of whose
$k$-consecutive sums are at most $\frac{k(mk+1)}{2}+1$.
\end{lemma}

\begin{proof}
Column $c$ of $A$ takes each value of $[m]$ once, so $\pi$ takes each value of the block
$\{m\beta(c)+1,\dots,m\beta(c)+m\}$ once, and these blocks partition $[mk]$. The $k$ consecutive positions starting
at $qk+r$ $(0\le r<k)$ are $qk+c$ for $c\ge r$ and $(q+1)k+c$ for $c<r$ (positions modulo $mk$, rows modulo $m$), one
in each column. Their sum is $m\sum_c\beta(c)+W_A(q,r)=\frac{mk(k-1)}{2}+W_A(q,r)\le\frac{mk(k-1)}{2}+\frac{k(m+1)}{2}+1
=\frac{k(mk+1)}{2}+1$.
\end{proof}

\begin{lemma}[base case]\label{lem:base}
Let $m=2h+1\ge3$ and $k=m-2$. Define $C:\Z_m\to[m]$ by $C(i)=m-i$ for $0\le i\le h$ and $C(i)=i-h$ for
$h<i\le 2h$, so that $C=(m,m-1,\dots,h+1,1,2,\dots,h)$. Then $A(q,c)=C(q+hc)$ is a good $m\times k$ array.
\end{lemma}

\begin{proof}
$C$ is a bijection onto $[m]$, so every column is a permutation of $[m]$. Since $2h=m-1\equiv-1\pmod m$ we have
$1\equiv-2h$, hence $A(q+1,c)=C(q+h(c-2))$. Therefore the entries of $W_A(q,r)$ are the values $C(q+hj)$ for
\[
  j\in\{-2,-1,\dots,r-3\}\cup\{r,r+1,\dots,k-1\}=\{-2,\dots,k-1\}\setminus\{r-2,r-1\}.
\]
The set $\{-2,\dots,k-1\}$ consists of $k+2=m$ consecutive integers and $\gcd(h,m)=1$, so the residues $q+hj$,
$-2\le j\le k-1$, run over all of $\Z_m$. With $y=q+h(r-2)$ we obtain
\[
  W_A(q,r)=\sum_{x\in\Z_m}C(x)-C(y)-C(y+h)=\frac{m(m+1)}{2}-\bigl(C(y)+C(y+h)\bigr).
\]
Now $C(y)+C(y+h)\ge m$ for every $y$: for $1\le y\le h$ it equals $(m-y)+y=m$; for $y=0$ it equals $m+(h+1)$; and for
$h<y\le 2h$ we have $y+h\equiv y-h-1\in\{0,\dots,h-1\}$, so it equals $(y-h)+(m-y+h+1)=m+1$. Hence
$W_A(q,r)\le\frac{m(m+1)}{2}-m=\frac{(m-2)(m+1)}{2}+1=\frac{k(m+1)}{2}+1$.
\end{proof}

\begin{lemma}[adding two columns]\label{lem:insert}
Let $A$ be a good $m\times k$ array and $0\le j<k$. Insert immediately before column $j$ two new columns, first a
copy $U$ of column $j$ and then its complement $\bar U=m+1-U$. The resulting $m\times(k+2)$ array $A'$ is good.
\end{lemma}

\begin{proof}
The new columns are permutations of $[m]$. In $A'$ the columns are, in order, the old columns $0,\dots,j-1$, then
$U,\bar U$, then the old columns $j,\dots,k-1$. Consider a staircase sum $W_{A'}(q,r')$, $0\le r'\le k+2$.
If $r'\le j$, both new columns are taken from row $q$; if $r'\ge j+2$, both are taken from row $q+1$. In these cases
they contribute $U+\bar U=m+1$, and the old columns contribute $W_A(q,r')$ or $W_A(q,r'-2)$ respectively, so
$W_{A'}(q,r')\le\frac{k(m+1)}2+1+(m+1)=\frac{(k+2)(m+1)}2+1$. If $r'=j+1$, then $U$ is taken from row $q+1$ and
$\bar U$ from row $q$, while the old columns contribute $W_A(q,j)$; thus
\[
  W_{A'}(q,j+1)=W_A(q,j)+A(q+1,j)+\bigl(m+1-A(q,j)\bigr)=W_A(q,j+1)+(m+1),
\]
and the same bound holds.
\end{proof}

\begin{proof}[Proof of Theorem~\ref{thm:construction}]
Start from the good $m\times(m-2)$ array of Lemma~\ref{lem:base} and apply Lemma~\ref{lem:insert}
$\frac{k-m+2}{2}$ times to obtain a good $m\times k$ array; then apply Lemma~\ref{lem:blocks}.
\end{proof}

For example, for $k=m=5$, inserting the pair before column $j=1$ of the base array for $m=5$ gives the columns
$54312$, $31254$, $35412$, $31254$, $25431$ (each listing rows $0$ to $4$), and $\beta(c)=4-c$ gives the
arrangement
\[
  25,18,13,8,2,\ 24,16,15,6,5,\ 23,17,14,7,4,\ 21,20,11,10,3,\ 22,19,12,9,1
\]
of $[25]$, whose $5$-consecutive sums are $63$, $65$ or $66$, while the average is $65$. In general the staircase
sums of the arrays in Theorem~\ref{thm:construction} take only the three values $\frac{k(m+1)}2+1$,
$\frac{k(m+1)}2$ and $\frac{k(m+1)}2-\frac{m-1}2$.

\section{Proofs of the corollaries}

\begin{proof}[Proof of Corollary~\ref{cor:mk}]
By Remark~\ref{rem:trivial}, $\msum(mk,k)\ge1$, and by Theorem~\ref{thm:construction} equality holds for
$m\le k+2$. For $m\ge k+4$ we have $g=\gcd(mk,k)=k$ and $mk>(k+3)k$, so Theorem~\ref{thm:lower}(a) gives
$\msum(mk,k)\ge2$, while $\msum(mk,k)\le\disc(mk,k)\le2$ for odd $k$ by~\cite{AFG} (see~\cite[Theorem~2.4(iv)]{HK}).
\end{proof}

\begin{proof}[Proof of Corollary~\ref{cor:clock}]
Here $a(n)=\frac{3(n+1)}2+\msum(n,3)$, and $f(n)=\lceil 3n/2\rceil+3$ equals $\frac{3(n+1)}2+2$ for odd $n$ and
$\frac{3(n+1)}2+\frac32$ for even $n$.

\emph{Upper bound.} For $n\ge6$, $\msum(n,3)\le\disc(n,3)\le2$~\cite{AFG} (see~\cite[Theorem~2.4(iii)]{HK}); for even $n$ this means
$\msum(n,3)\le\frac32$ because $\msum(n,3)\in\frac12+\Z$. So $a(n)\le f(n)$ for $n\ge6$. Schoenfield's explicit
arrangements~\cite{OEIS} give the same bound for all $n\ge3$. For $n=4$, every three neighbours omit one number, so
$a(4)=10-1=9=f(4)$.

\emph{Lower bound.} Let $n\ge5$ and $g=\gcd(n,3)$. If $n$ is odd and $3\nmid n$, Theorem~\ref{thm:lower}(a) gives
$\msum(n,3)\ge2$ for $n>6$. If $n=3m$ is odd, Corollary~\ref{cor:mk} gives $\msum(n,3)=2$ for $m\ge7$, that is for
$n\ge21$. If $n$ is even, Theorem~\ref{thm:lower}(b) gives $\msum(n,3)\ge\frac32$ unless $n=6$. Hence $a(n)=f(n)$ for all $n\ge4$ except possibly $n=5,6,9,15$.

\emph{The exceptions.} For odd $n$, $a(n)\ge\frac{3(n+1)}2+1=f(n)-1$ by Remark~\ref{rem:trivial}, and for $n=6$,
$a(6)\ge\lceil 21/2\rceil=11=f(6)-1$. The arrangements
\begin{gather*}
  (5,1,4,3,2),\qquad (6,1,4,5,2,3),\qquad (9,4,3,8,5,2,7,6,1),\\ (15,8,2,14,6,5,13,7,4,11,10,3,12,9,1)
\end{gather*}
have largest sums of three neighbours $10$, $11$, $16$ and $25$, respectively; the last two come from
Theorem~\ref{thm:construction} with $(k,m)=(3,3)$ and $(3,5)$.
\end{proof}

\section{Computations}

The arguments above need no computer. As a check we verified the following by computer; the programs are short
C and Python codes.
\begin{enumerate}
\item Lemma~\ref{lem:basic}, Lemma~\ref{lem:tv}, Proposition~\ref{prop:key}, Lemma~\ref{lem:pairs} and the final
inequality $n\le 2(T-n)+4g+(k-1)P$ hold for every circular arrangement of $[n]$ with $4\le n\le11$ $(k=3)$ and
$6\le n\le 11$ $(k=5)$, taking for $L$ the largest $k$-consecutive sum and the next two integers (about
$2.4\times10^7$ cases in all). A second, independently written program confirmed them for all $k\le9$ with
$n\le10$ and for $n=11,12$ with $k=3,5,7,9$ (Lemma~\ref{lem:pairs} and the final inequality for odd $k$ only).
\item The arrays of Theorem~\ref{thm:construction}, with several insertion positions and random block orders
$\beta$, give arrangements with largest $k$-consecutive sum exactly $\frac{k(mk+1)}2+1$ for all odd $3\le k\le61$
and odd $3\le m\le k+2$; the second program also tried every sequence of insertion positions for $k\le 11$ and
every block order $\beta$ for $k\le7$.
\item An exhaustive branch-and-bound search, which prunes with Lemmas~\ref{lem:basic}(ii) and~\ref{lem:tv} and was
checked against plain enumeration, reproduces all published values $a(3),\dots,a(38)$ and
gives $a(39)=62$, $a(40)=63$, $a(41)=65$, $a(42)=66$, $a(43)=68$ and $a(44)=69$, in agreement with
Corollary~\ref{cor:clock}. For $k=5$ it finds
arrangements with $\msum(\pi,5)=1$ for $n=15,25,35$, in agreement with Corollary~\ref{cor:mk}.
\end{enumerate}

\subsection*{Acknowledgement}
The arguments, the programs and the computations in this note were developed with the help of an AI assistant
(Claude, by Anthropic). The author has checked them and takes responsibility for the content.

\begin{thebibliography}{9}
\bibitem{AFG} R.~P.~Anstee, R.~Ferguson and J.~R.~Griggs, Permutations with low discrepancy consecutive $k$-sums,
\emph{J. Combin. Theory Ser. A} \textbf{100} (2002), no.~2, 302--321, \url{https://doi.org/10.1006/jcta.2002.3304}.
\bibitem{BMSM} M.~Bogdanovi\'c, Z.~Maksimovi\'c, A.~Simi\'c and J.~Milo\v{s}evi\'c, A mixed integer linear
programming formulation for low discrepancy consecutive $k$-sums permutation problem, \emph{Yugosl. J. Oper. Res.}
\textbf{27} (2017), no.~1, 125--132, \url{https://doi.org/10.2298/YJOR160104005B}.
\bibitem{BWM} Bundeswettbewerb Mathematik 2002, first round, problem~4.
\bibitem{Clark} D.~S.~Clark, A combinatorial theorem on circulant matrices, \emph{Amer. Math. Monthly} \textbf{92}
(1985), no.~10, \url{https://www.jstor.org/stable/2323225}.
\bibitem{Gardner} M.~Gardner, column in \emph{Isaac Asimov's Science Fiction Magazine}, August 1986.
\bibitem{HK} A.~Higashitani and K.~Kurimoto, Permutations with small maximal $k$-consecutive sums,
\emph{Discrete Math.} \textbf{342} (2019), 2467--2481; arXiv:1801.00416.
\bibitem{OEIS} OEIS Foundation Inc., The On-Line Encyclopedia of Integer Sequences, sequence A066385,
\url{https://oeis.org/A066385} (comments by J.~E.~Schoenfield, June and July 2019).
\bibitem{Rolfe} T.~Rolfe, Recurse around the clock, \emph{Mathematics and Computer Education} \textbf{21} (1987),
no.~2, 98--104.
\bibitem{RP} T.~J.~Rolfe and P.~W.~Purdom, An alternative problem for backtracking and bounding,
\emph{ACM SIGCSE Bulletin} \textbf{36} (2004), no.~4, 83--84.
\end{thebibliography}

\end{document}
